\subsection{Voorbereiding data}
\label{subsec:voorbereiding-data}

De in sectie~\ref{subsubsec:datasets} besproken corpora zijn niet geschikt voor dit onderzoek, omdat ze niet de vereiste e-mailheaders bevatten. Daarom wordt geen gebruikgemaakt van MeAJOR \autocite{Mendes2025}, Enron \autocite{klimt2004} of Nazario \autocite{Nazario2025}, maar van de volgende twee datasets:
\begin{enumerate}
    \item \textbf{Phishing\_pot.} Een verzameling phishingmails die actief onderhouden wordt door \textcite{PhishingPot};
    \item \textbf{SpamAssassin Public Corpus.} Een verzameling e-mailberichten uit de vroege jaren 2000, met zowel ongewenste als legitieme e-mails \footfullcite{SpamAssassinCorpus}. Voor dit onderzoek wordt echter uitsluitend gebruikmaakt van de partities \\ \emph{easy\_ham} en \emph{hard\_ham}. De eerste bevat legitieme e-mails die relatief eenvoudig te detecteren zijn. De tweede bevat legitieme e-mails die meer gelijkenissen vertonen met frauduleuze berichten.
\end{enumerate}

\paragraph{Beperkingen van de data.}
Het SpamAssassin-project raadt zelf af om hun legitieme e-mailset te gebruiken als die ouder is dan 18 maanden is, omdat legitieme e-mail met de tijd verandert \footfullcite{SpamAssassinCorpusPolicy}.
Het gebruik van deze data blijft echter verdedigbaar, omdat het doel een risicoscore is op basis van structurele kenmerken, zoals headerconsistentie, URL-opbouw of hyperlinks. Deze kenmerken evolueren naar verwachting trager dan woordgebruik.

\paragraph{Voorbereiding.} 
Elke e-mail wordt ingelezen en omgezet naar een verwerkbare vorm. Daarbij worden headers en body gedecodeerd waar nodig, zoals \emph{encoded-words} in headers (RFC~2047 \autocite{Moore1996}) en body-delen die als base64 zijn gecodeerd. 

Vervolgens wordt de dataset opgesplitst in een trainingsset van 70\% en een testset van 30\%. Hierbij blijft de verhouding tussen phishingmails en legitieme e-mails in beide sets behouden. Een vaste seed zorgt ervoor dat de opsplitsing bij herhaaldelijke uitvoering reproduceerbaar is.

Uit elk bericht worden vervolgens de kenmerken berekend. Deze kenmerken, samen met de indeling van de berichten in trainings- en testset, worden opgeslagen. Zo hoeven de berichten niet bij elk experiment opnieuw te worden verwerkt, en worden alle modelconfiguraties op dezelfde gegevens getraind.

De testset wordt pas gebruikt wanneer de parameters van het model op punt staan na voldoende trainingsiteraties.

